\scanpage{79}
\mspage{23}

Essayons de le prouver, en calquant la dém. dans le cas $C=C_\infty$ (\uline{XI}, th.~2). On peut supposer $S$ \uline{connexe}, en \deleted{supposant} utilisant $\mathrm{L}\,\unsure{\mathrm{a}|0}$ (i.e. $W\subset W_{\unsure{0|o}}$) et en admettant le
% slanted
\margin{ne sert à rien !}

$\blacklozenge$ \uline{Lemme}. Soit $f:X\to Y$ dans $\mathrm{Cat}$. Alors $f\in W$ ssi $\pi_0(f)$ est bijectif, et pour \unsure{tte} composante \unsure{conn.|connexe} $X_i$ de $X$, \unsure{\{|$f$} si $Y_i$ est la comp. conn. corr. de $Y$, $f_i:X_i\to Y_i$ induit par $f$ est dans $W$.

Utilisons alors l'existence de
\[ S'\longrightarrow S \]
Hot-fibration avec $S'$ Hot-asphérique. \unsure{OPS} = $S'\unsure{\text{---}|\to}S$ est parfait. Supposons qu'on sache qu'un chgt. de base par de telles fibrations parfaites transforme \deleted{chgt. de} $W$-équivalence en itou. (Ça doit être vrai pour les $W_i$, mais on a vu que c'était faux pour l'équivalence \unsure{ho-homologique}, cf I p.~87 ff) Alors
% slanted
\margin{\uline{NB} Il suffit d'avoir une Hot-fibration $S'\unsure{\text{---}|\to}S_0\subset S$, où $S_0$ est la comp. conn. de $S$ contenant l'élément $s\in\operatorname{Ob}S$ envisagé.}
$X'\xrightarrow{f'}Y'$ sur $S'$ est une $W$-équivalence. Or pour $s'\in S'$, comme $\{s'\}\longrightarrow S'$ est une Hot-équivalence ($S'$ Hot-asphérique), et $X',Y'$ étant Hot-fibrés sur $S$, il s'ensuit que $X'_{s'}\to X',\ Y'_{s'}\unsure{\text{---}|\to}Y'$ sont des Hot-équivalences, donc aussi $X_{s'}\overset{f'_{s'}}{\text{---}}Y_{s'}$, comme on voit sur le carré comm.
\diagram{\begin{tikzcd}[ampersand replacement=\&]
X'_{s'}\arrow[r,"f'_{s'}"]\arrow[d,"\sim"']\& Y'_{s'}\arrow[d,"\sim"]\\
X'\arrow[r,"f'"']\&Y'
\end{tikzcd}}

% archive sheet number: 78
